% Part: second-order-logic
% Chapter: sol-and-set-theory
% Section: power-of-continuum

\documentclass[../../../include/open-logic-section]{subfiles}

\begin{document}

\olfileid{sol}{set}{pow}

\olsection{قوة المتصل}

\begin{explain}
التكميم في الرتبة الثانية يتناول مجموعات !!{domain} الجزئية،
ولا يتناول مجموعات تلك المجموعات الجزئية من !!{domain}.
ولو أردنا هذا التكميم مباشرة لاحتجنا إلى منطق \emph{الرتبة الثالثة}.
فمثلًا، مبرهنة كانتور تنفي وجود دالة !!{injective} من مجموعة
قوى مجموعة إلى المجموعة نفسها. وقد نحاول أن نقول:
«لكل مجموعة~$\OLId{latin-looped}{X}$ ولكل مجموعة~$\OLId{latin-looped}{P}$، إذا كانت~$\OLId{latin-looped}{P}$ مجموعة
قوى~$\OLId{latin-looped}{X}$، فليس~$\cardle{\OLId{latin-looped}{P}}{\OLId{latin-looped}{X}}$».
لكن كون~$\OLId{latin-looped}{P}$ مجموعة قوى~$\OLId{latin-looped}{X}$ معناه أن !!{element}s~$\OLId{latin-looped}{P}$
هي مجموعات~$\OLId{latin-looped}{X}$ الجزئية كلها ولا غيرها، فنحتاج إلى عبارة
من نحو~$\lforall[\OLId{latin-looped}{Y}][(\OLId{latin-looped}{P}(\OLId{latin-looped}{Y}) \liff \OLId{latin-looped}{Y} \subseteq \OLId{latin-looped}{X})]$.
وموضع الإشكال هو~$\OLId{latin-looped}{P}(\OLId{latin-looped}{Y})$؛ إذ ليست هذه !!a{formula} من الرتبة
الثانية: متغير العلاقة الأحادي الموضع، كـ$\OLId{latin-looped}{P}$، لا يأخذ في موضع
وسيطه إلا حدًا.

ومع ذلك يمكن \emph{محاكاة} التكميم على مجموعات المجموعات
إذا اتسع المجال لذلك. نستخدم علاقة ثنائية الموضع~$\OLId{latin-looped}{R}$،
تربط !!{element}s من !!{domain} بـ!!{element}s من !!{domain}.
فمتى أعطينا~$\OLId{latin-looped}{R}$ جمعنا !!{element}s التي يرتبط بها عنصر~$\OLId{latin-ordinary}{x}$
بالعلاقة~$\OLId{latin-looped}{R}$. وتكون~$\Setabs{\OLId{latin-ordinary}{y} \in \Domain{\OLId{latin-looped}{M}}}{\OLId{latin-looped}{R}(\OLId{latin-ordinary}{x},\OLId{latin-ordinary}{y})}$
المجموعة التي نرمز إليها بـ$\OLId{latin-ordinary}{x}$.
وعكسًا، لتكن~$\OLId{latin-looped}{Z} \subseteq \Pow{\Domain{\OLId{latin-looped}{M}}}$ مجموعة من
مجموعات~$\Domain{\OLId{latin-looped}{M}}$ الجزئية. فإن كان في~$\Domain{\OLId{latin-looped}{M}}$ من
!!{element}s ما لا يقل عددًا عن مجموعات~$\OLId{latin-looped}{Z}$، وجدت
علاقة~$\OLId{latin-looped}{R} \subseteq \Domain{\OLId{latin-looped}{M}}^\OLMathNumeral{2}$ يكون بها لكل مجموعة رمز:
بوساطة~$\OLId{latin-looped}{R}$ يرمز بعض~$\OLId{latin-ordinary}{x}$ إلى كل~$\OLId{latin-looped}{Y} \in \OLId{latin-looped}{Z}$.
\end{explain}

\begin{defn}
متى كان~$\OLId{latin-looped}{R} \subseteq \Domain{\OLId{latin-looped}{M}}^\OLMathNumeral{2}$ قلنا إن
\emph{$\OLId{latin-ordinary}{x}$ يرمز بالعلاقة~$\OLId{latin-looped}{R}$} إلى
$\Setabs{\OLId{latin-ordinary}{y} \in \Domain{\OLId{latin-looped}{M}}}{\OLId{latin-looped}{R}(\OLId{latin-ordinary}{x}, \OLId{latin-ordinary}{y})}$.
\end{defn}

إذا كان !!a{element}~$\OLId{latin-ordinary}{x} \in \Domain{\OLId{latin-looped}{M}}$ يرمز
بالعلاقة~$\OLId{latin-looped}{R}$ إلى مجموعة~$\OLId{latin-looped}{Z} \subseteq \Domain{\OLId{latin-looped}{M}}$،
كانت المجموعة~$\OLId{latin-looped}{Y} \subseteq \Domain{\OLId{latin-looped}{M}}$ ترمز إلى جملة من المجموعات:
ما ترمز إليه !!{element}s~$\OLId{latin-looped}{Y}$.
وبهذا قد تكون~$\OLId{latin-looped}{Y}$ رمزًا، بالعلاقة~$\OLId{latin-looped}{R}$، لـ$\Pow{\OLId{latin-looped}{X}}$.
وشرط ذلك أن يكون لكل~$\OLId{latin-looped}{Z} \subseteq \OLId{latin-looped}{X}$ عنصر~$\OLId{latin-ordinary}{x} \in \OLId{latin-looped}{Y}$
يرمز بالعلاقة~$\OLId{latin-looped}{R}$ إلى~$\OLId{latin-looped}{Z}$، وأن يكون كل~$\OLId{latin-ordinary}{x} \in \OLId{latin-looped}{Y}$
رمزًا بالعلاقة~$\OLId{latin-looped}{R}$ لمجموعة~$\OLId{latin-looped}{Z} \subseteq \OLId{latin-looped}{X}$.

\begin{prop}
!!{formula}
  \[
  \fn{Codes}(\OLId{latin-ordinary}{x}, \OLId{latin-looped}{R}, \OLId{latin-looped}{Z}) \ident \lforall[\OLId{latin-ordinary}{y}][(\OLId{latin-looped}{Z}(\OLId{latin-ordinary}{y}) \liff
    \OLId{latin-looped}{R}(\OLId{latin-ordinary}{x}, \OLId{latin-ordinary}{y}))]
  \]
تقرر أن~$\OLId{latin-ordinary}{s}(\OLId{latin-ordinary}{x})$ يرمز بالعلاقة~$\OLId{latin-ordinary}{s}(\OLId{latin-looped}{R})$ إلى~$\OLId{latin-ordinary}{s}(\OLId{latin-looped}{Z})$.
وأما !!{formula}
\begin{multline*}
  \fn{Pow}(\OLId{latin-looped}{Y}, \OLId{latin-looped}{R}, \OLId{latin-looped}{X}) \ident \\
  \lforall[\OLId{latin-looped}{Z}][(\OLId{latin-looped}{Z} \subseteq \OLId{latin-looped}{X} \lif \lexists[\OLId{latin-ordinary}{x}][(\OLId{latin-looped}{Y}(\OLId{latin-ordinary}{x}) \land
    \fn{Codes}(\OLId{latin-ordinary}{x}, \OLId{latin-looped}{R}, \OLId{latin-looped}{Z}))])] \land {} \\ \lforall[\OLId{latin-ordinary}{x}][(\OLId{latin-looped}{Y}(\OLId{latin-ordinary}{x}) \lif
  \lforall[\OLId{latin-looped}{Z}][(\fn{Codes}(\OLId{latin-ordinary}{x}, \OLId{latin-looped}{R}, \OLId{latin-looped}{Z}) \lif \OLId{latin-looped}{Z} \subseteq \OLId{latin-looped}{X})])]
\end{multline*}
فتقرر أن~$\OLId{latin-ordinary}{s}(\OLId{latin-looped}{Y})$ ترمز بالعلاقة~$\OLId{latin-ordinary}{s}(\OLId{latin-looped}{R})$ إلى مجموعة قوى~$\OLId{latin-ordinary}{s}(\OLId{latin-looped}{X})$؛
أي إن عناصر~$\OLId{latin-ordinary}{s}(\OLId{latin-looped}{Y})$ ترمز بالعلاقة~$\OLId{latin-ordinary}{s}(\OLId{latin-looped}{R})$ إلى مجموعات~$\OLId{latin-ordinary}{s}(\OLId{latin-looped}{X})$
الجزئية كلها دون غيرها.
\end{prop}

\begin{explain}
بهذا الترميز نعبّر عن قضايا مجموعة القوى، فنجعل التكميم على
رموز المجموعات الجزئية عوضًا عن المجموعات نفسها.
ومن ذلك أن نصوغ مبرهنة كانتور بنفي وجود دالة !!{injective}
من أي مجموعة~$\OLId{latin-looped}{Y}$ ترمز بالعلاقة~$\OLId{latin-looped}{R}$ إلى مجموعة قوى~$\OLId{latin-looped}{X}$، إلى~$\OLId{latin-looped}{X}$ نفسها.
\end{explain}

\begin{prop}
الجملة التالية صحيحة منطقيًا:
\begin{multline*}
  \lforall[\OLId{latin-looped}{X}][\lforall[\OLId{latin-looped}{Y}][\lforall[\OLId{latin-looped}{R}][(\fn{Pow}(\OLId{latin-looped}{Y}, \OLId{latin-looped}{R}, \OLId{latin-looped}{X}) \lif \\
      \lnot \lexists[\OLId{latin-ordinary}{u}][(\lforall[\OLId{latin-ordinary}{x}][\lforall[\OLId{latin-ordinary}{y}][((\OLId{latin-looped}{Y}(\OLId{latin-ordinary}{x}) \land \OLId{latin-looped}{Y}(\OLId{latin-ordinary}{y}) \land
            \eq[\OLId{latin-ordinary}{u}(\OLId{latin-ordinary}{x})][\OLId{latin-ordinary}{u}(\OLId{latin-ordinary}{y})]) \lif
            \eq[\OLId{latin-ordinary}{x}][\OLId{latin-ordinary}{y}])]] \land {}\\
        \lforall[\OLId{latin-ordinary}{x}][(\OLId{latin-looped}{Y}(\OLId{latin-ordinary}{x}) \lif \OLId{latin-looped}{X}(\OLId{latin-ordinary}{u}(\OLId{latin-ordinary}{x})))])])]]]
\end{multline*}
\end{prop}

\begin{explain}
مجموعة قوى مجموعة !!a{denumerable} هي !!{nonenumerable}؛
فأصليتها تزيد على أصلية أي مجموعة !!{denumerable}، أي
$\aleph_\OLMathNumeral{0}$. وحجم~$\Pow{\Nat}$ هو «قوة المتصل»، إذ يساوي
حجم مجموعة نقاط المستقيم الحقيقي~$\Real$.
فإن اتسع !!{domain} لترميز مجموعة قوى مجموعة !!a{denumerable}،
  أمكن أن نصف مجموعة بأنها من حجم المتصل بقولنا إنها تكافئ
  عدديًا مجموعة~$\OLId{latin-looped}{Y}$ ترمز مجموعة قوى مجموعة~$\OLId{latin-looped}{X}$
  من الحجم~$\aleph_\OLMathNumeral{0}$، بشرط أن يكون لكل مجموعة جزئية من~$\OLId{latin-looped}{X}$
  رمز واحد لا غير في~$\OLId{latin-looped}{Y}$؛ وهو شرط التفرّد المثبت في
  \(\fn{Cont}\) أدناه. فأما الترميز غير المتفرّد فلا يفيد إلا
  حدًا أدنى لحجم~$\OLId{latin-looped}{Y}$. وإن ضاق المجال عن ذلك، فلم توجد
فيه مجموعة جزئية مكافئة عدديًا لـ$\Real$، امتنع وجود علاقة
  ترمز~$\Pow{\OLId{latin-looped}{X}}$ أيضًا.
% تصحيح مصدر (0340-I1، 0340-I2): قُيد الحقن بالمجموعة Y وصُرّح بشرط تفرد رموز مجموعة القوى.
\end{explain}

\begin{prop}
إذا كان~$\cardle{\Real}{\Domain{\OLId{latin-looped}{M}}}$ كانت !!{formula}
\begin{multline*}
\fn{Cont}(\OLId{latin-looped}{Y}) \ident 
\lexists[\OLId{latin-looped}{X}][\lexists[\OLId{latin-looped}{R}][((\fn{Aleph}_\OLMathNumeral{0}(\OLId{latin-looped}{X}) \land
      \fn{Pow}(\OLId{latin-looped}{Y}, \OLId{latin-looped}{R}, \OLId{latin-looped}{X})) \land \\ 
      \lforall[\OLId{latin-ordinary}{x}][\lforall[\OLId{latin-ordinary}{y}][((\OLId{latin-looped}{Y}(\OLId{latin-ordinary}{x}) \land {} 
      \OLId{latin-looped}{Y}(\OLId{latin-ordinary}{y}) \land \lforall[\OLId{latin-ordinary}{z}][\OLId{latin-looped}{R}(\OLId{latin-ordinary}{x},\OLId{latin-ordinary}{z}) \liff \OLId{latin-looped}{R}(\OLId{latin-ordinary}{y},\OLId{latin-ordinary}{z})]) \lif \eq[\OLId{latin-ordinary}{x}][\OLId{latin-ordinary}{y}])]])]]
\end{multline*}
تعبيرًا عن~$\cardeq{\OLId{latin-ordinary}{s}(\OLId{latin-looped}{Y})}{\Real}$.
\end{prop}

\begin{proof}
تقرر~$\fn{Pow}(\OLId{latin-looped}{Y}, \OLId{latin-looped}{R}, \OLId{latin-looped}{X})$ أن~$\OLId{latin-ordinary}{s}(\OLId{latin-looped}{Y})$ ترمز
بالعلاقة~$\OLId{latin-ordinary}{s}(\OLId{latin-looped}{R})$ إلى مجموعة قوى~$\OLId{latin-ordinary}{s}(\OLId{latin-looped}{X})$؛
وتقرر~$\fn{Aleph}_\OLMathNumeral{0}(\OLId{latin-looped}{X})$ أن المجموعة الأصلية قابلة للتعداد.
فـ$\OLId{latin-ordinary}{s}(\OLId{latin-looped}{Y})$ لا يقل حجمها عن قوة المتصل، وقد يزيد عليها إذا
اجتمعت عناصر عدة من~$\OLId{latin-ordinary}{s}(\OLId{latin-looped}{Y})$ ترمز إلى الجزء نفسه من~$\OLId{latin-looped}{X}$.
وإنما يمنع المقترن الأخير هذا الاحتمال؛ إذ يشترط أن يكون الربط
بالعلاقة~$\OLId{latin-ordinary}{s}(\OLId{latin-looped}{R})$ بين عناصر~$\OLId{latin-ordinary}{s}(\OLId{latin-looped}{Y})$ ومجموعات~$\OLId{latin-ordinary}{s}(\OLId{latin-looped}{X})$ الجزئية
!!{injective}.
\end{proof}

\begin{prop}
يتحقق~$\cardeq{\Domain{\OLId{latin-looped}{M}}}{\Real}$ إذا وفقط إذا
\begin{multline*}
  \Sat{\OLId{latin-looped}{M}}{\lexists[\OLId{latin-looped}{X}][\lexists[\OLId{latin-looped}{Y}][\lexists[\OLId{latin-looped}{R}][(\fn{Aleph}_\OLMathNumeral{0}(\OLId{latin-looped}{X}) \land \fn{Pow}(\OLId{latin-looped}{Y}, \OLId{latin-looped}{R}, \OLId{latin-looped}{X}) \land {}\\
      \lforall[\OLId{latin-ordinary}{x}][\lforall[\OLId{latin-ordinary}{y}][((\OLId{latin-looped}{Y}(\OLId{latin-ordinary}{x}) \land \OLId{latin-looped}{Y}(\OLId{latin-ordinary}{y}) \land
        \lforall[\OLId{latin-ordinary}{z}][\OLId{latin-looped}{R}(\OLId{latin-ordinary}{x},\OLId{latin-ordinary}{z}) \liff \OLId{latin-looped}{R}(\OLId{latin-ordinary}{y},\OLId{latin-ordinary}{z})]) \lif \eq[\OLId{latin-ordinary}{x}][\OLId{latin-ordinary}{y}])]] \land {}\\
      \lexists[\OLId{latin-ordinary}{u}][(\lforall[\OLId{latin-ordinary}{x}][\lforall[\OLId{latin-ordinary}{y}][(\eq[\OLId{latin-ordinary}{u}(\OLId{latin-ordinary}{x})][\OLId{latin-ordinary}{u}(\OLId{latin-ordinary}{y})] \lif \eq[\OLId{latin-ordinary}{x}][\OLId{latin-ordinary}{y}])]] \land {}
        \\\lforall[\OLId{latin-ordinary}{x}][\OLId{latin-looped}{Y}(\OLId{latin-ordinary}{u}(\OLId{latin-ordinary}{x}))] \land {}
        \\\lforall[\OLId{latin-ordinary}{y}][(\OLId{latin-looped}{Y}(\OLId{latin-ordinary}{y}) \lif \lexists[\OLId{latin-ordinary}{x}][\eq[\OLId{latin-ordinary}{y}][\OLId{latin-ordinary}{u}(\OLId{latin-ordinary}{x})]])])])]]]}. 
\end{multline*}
\end{prop}

\begin{explain}
مؤدى فرضية المتصل أن حجم المتصل هو أول أصلية
!!{nonenumerable}؛ أي إن حجم~$\Pow{\Nat}$ هو~$\aleph_\OLMathNumeral{1}$.
\end{explain}

\begin{prop}
صدق فرضية المتصل يكافئ الصحة المنطقية للجملة
\[\fn{CH} \ident
\lforall[\OLId{latin-looped}{X}][(\fn{Aleph}_\OLMathNumeral{1}(\OLId{latin-looped}{X}) \liff \fn{Cont}(\OLId{latin-looped}{X}))]\]
\end{prop}

ولا يلزم من هذا أن تكون~$\lnot \fn{CH}$ صحيحة منطقيًا إذا
وفقط إذا كذبت فرضية المتصل. فالمجال !!a{enumerable} لا يشتمل
على مجموعة جزئية حجمها~$\aleph_\OLMathNumeral{1}$، ولا على مجموعة جزئية
من حجم المتصل. لذلك تصدق~$\fn{CH}$ دائمًا في مجال
!!a{enumerable}. غير أن لنفي الفرضية جملة أخرى تكون
صحيحة منطقيًا إذا وفقط إذا كذبت الفرضية:

\begin{prop}
كذب فرضية المتصل يكافئ الصحة المنطقية للجملة
\[\fn{NCH} \ident
\lforall[\OLId{latin-looped}{X}][(\fn{Cont}(\OLId{latin-looped}{X}) \lif \lexists[\OLId{latin-looped}{Y}][(\OLId{latin-looped}{Y} \subseteq \OLId{latin-looped}{X} \land \lnot
    \fn{Count}(\OLId{latin-looped}{Y}) \land \lnot \cardeq{\OLId{latin-looped}{X}}{\OLId{latin-looped}{Y}})])]\]
\end{prop}

\end{document}
